\chapter{Future Work}
\label{chap:futureWork}

Using \glspl{gpzkp} for showing ballot validity is still a relatively new idea. 
Many of the components used in our implementation can be optimized further or interchanged with others.

As mentioned in \cref{chap:software,chap:results}, using the default compilation option \verb|--O1| for Circom circuits in combination with Groth16 as the backend prover is not ideal.
Compiling with the option \verb|--O1| creates very versatile \glspl{r1cs} which also allow other backend provers such as PLONK and FFLONK \cite{gabizon_plonk_2019,gabizon_fflonk_2021}.
However, for Groth16 the compilation option \verb|--O2| is preferrable since it leads to \glspl{r1cs} with lower total constraint counts.
Therefore, recompiling the ballot validity circuits $\mathfrak{C}_C^{Enc}$, $\mathfrak{C}_C^{Vot}$ and $\mathfrak{C}_C^{Ballot\text{-}Validity}$ with the option \verb|--O2| should improve our results.

The implementations of $\mathfrak{C}_C^{Vot}$ are fairly simple and contribute only little to the total number of constraints generated by $\mathfrak{C}_C^{Ballot\text{-}Validity}$.
Therefore, the potential for optimization in terms of constraint count and subsequently proving time as well as $CRS$ size is comparatively small.

However, as shown in our work, optimizing the circuit $\mathfrak{C}_C^{Enc}$ can reduce the constraints, proving time and $CRS$ size significantly.
In the future one could try to optimize our version of the circuit that uses \gls{ppeeg} encryption schemes and Twisted Edwards curves further.
A potential improvement to the subcircuit $\mathfrak{C}_{TE_{a,d}(\mathbb{F})\_PPFExp}$ which computes group exponentiations $P^k$ could be to change the binary representation of the exponent $k$ to a represenation with a different base.
Additionally, providing the exponent $k$ in the representation to the chosen base dircetly would remove the need for computing this representation inside the circuit.
For instance, providing $k$ in a representation to base $5$ could reduce the constraint count by up to $35\%$ as outlined in \cref{section:appendix:ppeegArbitraryBase}.

Another way to reduce the constraint count of $\mathfrak{C}_C^{Enc}$ could be to employ a different type of elliptic curve.
Note though, that directly exchanging the Twisted Edwards curve in $\mathfrak{C}_{TE_{a,d}(\mathbb{F})\_PPFExp}$ should not yield any reductions since the circuit $\mathfrak{C}_{TE_{a,d}(\mathbb{F})\_\odot}$ implementing the Twisted Edwards group law uses only $7$ constraints.
This is the least out of any of the commonly used ellitic curve types Short Weierstrass, Montgomery and Twisted Edwards.
Nevertheless, one might be able to use pseudooperations only involving a part of the coordinates to reduce the constraints further, similar to how those were employed for Montgomery curves in the Montgomery Ladder.
% We explore this idea in more detail in \cref{todo}\todo{Make in appendix} and show why straightforward integrations of the pseudooperations are not possible or use too many constraints.

Furthermore, the elliptic curve \verb|BN254| we used as the basis for the \gls{g16snark} could be exchanged for a different one.
Potential candidates are curves based on the BLS12-381 curve \cite{barreto_constructing_2002, sean_bowe_bls12-381_2017} which offer a slightly higher security level than \verb|BN254|
This was already suggested previously by Huber et. al. \cite{huber_zk-snarks_2025}.

In addition to that, there is potential for future research outside of the implementation of the circuits themself.
In our implementation we used Circom and the backend prover snarkjs to define circuits and compute \glspl{g16snark} over them while Huber et. al. employed libsnark for this purpose.
However, snarkjs also supports other \glspl{zps} like PLONK and a variant FFLONK \cite{gabizon_plonk_2019, gabizon_fflonk_2021}.
Compared to \gls{g16snark} these \glspl{zps} have the advantage of not needing a circuit specific $CRS$.
However, computing proofs is comparatively inefficient \cite{gabizon_plonk_2019,gabizon_fflonk_2021}. % as outlined in \cref{todo} \todo{Make in appendix}.
% However, our implementation is not designed for these \glspl{zps}
Nevertheless, an in-depth comparison between an instantiation of our implementation with PLONK and FFLONK and the \gls{g16snark} instantiation might be interesting.

Furthermore, Circom is not domain specific and can therefore be used in combination with many different backend provers such as arkworks \cite{yuncong_hu_arkworks_2025}.
arkworks also supports even more different \glspl{zps} such as marlin and gemini \cite{chiesa_marlin_2019,bootle_gemini_2022}.
Another interesting choice might be the icicle-snark \cite{ingonyama_ingonyama-zkicicle-snark_2025} which supports Groth16 and is based on icicle \cite{ingonyama_icicle_2024, ingonyama_introducing_2023}, a hardware acceleration library for compute intensive cryptography.
For their implementation of \gls{g16snark}, the authors reported speedups of $20x$ to $100x$ compared to snarkjs.
Moreover, icicle-snark uses the same \verb|witness.wtns| and \verb|key.zkey| files to create the proofs as snarkjs (\cref{section:software:snarkjs}).
Therefore, testing our implementation with the icicle-snark as the backend prover should be relatively straightforward.

Additionally, there are many more circuit languages such as ZoKrates and gnark \cite{zokrates_zokrateszokrates_2025, consensys_consensysgnark_2025}.
Therefore, a comparison of the performance of our ballot validity proofs using our approach to an implementation in a different circuit langage or with a different backend prover and \gls{zps} might yield interesting results.

Finally, another interesting topic of research would be to compare the performance of \glspl{gpzkp} for ensuring ballot validity to that of specialized \glspl{zkp} for the election types where this is possible.
For instance, for the election type Single-Vote, there are voting schemes using ElGamal encryption and making use of Chaum-Pedersen proofs \cite{cramer_secure_1997}.
Comparing this to our implementation could be interesting to quantify the how much faster the specialized approach is.
This information could allow practitioners to better make an informed choice between the two different approaches to voting schemes.
The results of such a comparison could also be used to evaluate whether optimizing specialized \glspl{zkp} for ballot validity further is worthwhile or whether the difference to the \glspl{gpzkp} based approach is too small to justify this.


